\chapter{Desempeño en simulación}\label{chap:desempeno}

El capítulo~\ref{chap:controlador} diseñó el controlador PII sobre el modelo linealizado en $-4$~cm. Este capítulo responde a la pregunta que sigue: qué hace ese controlador cuando se conecta a la planta completa, con su fuerza electromagnética no lineal, su zona muerta y su fricción estática. Se delimita primero el intervalo en que el actuador permite operar; después se prueba el lazo en regulación, primero sin y luego con atascamiento-deslizamiento, y se aísla mediante el modelo lineal la causa de lo que aparece. Se evalúa a continuación el seguimiento de referencias senoidales y trapezoidales y, por último, el comportamiento en la configuración estable. En todo el capítulo se \emph{describe} lo que se observa; el capítulo~\ref{chap:doble_integral} explica por qué ocurre.

\section{Intervalo de operación}
\label{sec:intervalo}

El actuador está acotado a $u\in[0;\,3{,}5]$~V, y la sección~\ref{sec:actuador_estatico} del capítulo~\ref{chap:modelo} muestra que esa cota deja al imán sostenido pero inmóvil en $-4$~cm, el punto de diseño. Para delimitar en qué posiciones sí puede desplazarse, se simularon 7503 escalones con distintas posiciones iniciales y finales, con el PII y con el PI, en presencia de atascamiento-deslizamiento. La tabla~\ref{tab:intervalo} resume el resultado con el PII.

\begin{table}[ht]
\centering
\caption{Referencias alcanzables con $u\le3{,}5$~V desde distintas posiciones iniciales (PII, con atascamiento-deslizamiento).}
\label{tab:intervalo}
\begin{tabular}{|c|c|}
\hline
Posición inicial [cm] & Referencias alcanzables [cm] \\ \hline
$-4{,}0$ & ninguna: el imán se sostiene pero no se mueve \\ \hline
$-3{,}5$ & de $-3{,}6$ a $-2{,}0$ \\ \hline
$-3{,}0$ & de $-3{,}5$ a $-1{,}5$ \\ \hline
$-2{,}0$ & de $-3{,}3$ a $-0{,}5$ \\ \hline
$-1{,}0$ & de $-2{,}5$ a $+0{,}5$ \\ \hline
\end{tabular}
\end{table}

Los cambios de referencia son viables a partir de unos $-3{,}6$~cm. Por ello, las pruebas de este capítulo se realizan en el intervalo de $-3$ a $-2$~cm, con un escalón de $-2$ a $-3$~cm que deja margen respecto a ambos límites. Los márgenes de estabilidad del PII en ese intervalo se presentaron en la sección~\ref{sec:robustez_pii}.

\section{Simulación no lineal sin atascamiento-deslizamiento}

La primera prueba establece la referencia ideal. El controlador PII se evaluó sobre el modelo no lineal del levitador con la fuerza electromagnética no lineal y la fricción viscosa únicamente,
\begin{equation}
m\ddot{x} = \frac{u}{b\,(a-x)^4} - m g - m c_1 \dot{x} ,
\label{eq:nl_sin_ad}
\end{equation}
es decir, sin la fricción estática de Karnopp ni la zona muerta del actuador que las secciones~\ref{sec:friccion_modelo} y~\ref{sec:zm_modelo} incorporan a la ecuación de movimiento~\eqref{eq:mec}. Las condiciones son las de la prueba de regulación de la sección~\ref{sec:regulacion}: escalón de referencia de $-2$ a $-3$~cm en $t=10$~s, señal de control acotada a $u\in[0;\,3{,}5]$~V y controlador inicializado en el estado estacionario correspondiente a $x_1=-2$~cm. Para facilitar la comparación, cada figura superpone el resultado con atascamiento-deslizamiento (AD), que se presenta en esa misma sección.

\begin{figure}[ht]
    \centering
    % Estilo común de las figuras sin/con atascamiento-deslizamiento (pgfplots).
\pgfplotsset{
  sinad/.style={
    width=0.92\textwidth, height=5.2cm, xmin=0, xmax=40,
    xlabel={tiempo [\si{s}]},
    grid=major, grid style={gray!25},
    tick label style={font=\small}, label style={font=\small}, legend style={font=\small},
    /pgf/number format/use comma,
    scaled ticks=false,
    y tick label style={/pgf/number format/fixed},
    table/col sep=comma,
    every axis plot/.append style={line join=round},
  },
  detalle/.style={height=4cm, axis background/.style={fill=ctl2!4}},
  sinAD/.style={tray, line width=1.1pt},
  conAD/.style={ctl2, line width=0.7pt},
}
\definecolor{tray}{RGB}{31,94,166}
\definecolor{refc}{RGB}{40,40,40}
\definecolor{ctl2}{RGB}{196,96,40}

\begin{tikzpicture}
\begin{axis}[sinad, name=principal, xlabel={}, ylabel={posición [\si{cm}]}, ymin=-3.35, ymax=-1.85,
             legend style={at={(0.99,0.97)}, anchor=north east}, legend cell align=left]
  \addplot[conAD] table[x=t, y=x_ad] {images/sin_ad_results/tikz/comparacion.csv};
  \addlegendentry{con AD}
  \addplot[sinAD] table[x=t, y=x] {images/sin_ad_results/tikz/comparacion.csv};
  \addlegendentry{sin AD}
  \addplot[refc, densely dashed, line width=0.8pt] table[x=t, y=r] {images/sin_ad_results/tikz/comparacion.csv};
  \addlegendentry{referencia}
  \fill[ctl2, opacity=0.15] (axis cs:9.8,-3.35) rectangle (axis cs:12,-1.85);
\end{axis}
\begin{axis}[sinad, detalle, at={(principal.south west)}, anchor=north west, yshift=-0.45cm,
             xmin=9.8, xmax=12, ymin=-3.3, ymax=-1.9, ylabel={detalle [\si{cm}]}]
  \addplot[conAD, line width=1.6pt] table[x=t, y=x_ad] {images/sin_ad_results/tikz/detalle.csv};
  \addplot[sinAD, line width=0.8pt] table[x=t, y=x] {images/sin_ad_results/tikz/detalle.csv};
  \addplot[refc, densely dashed, line width=0.8pt] table[x=t, y=r] {images/sin_ad_results/tikz/detalle.csv};
\end{axis}
\end{tikzpicture}

    \caption[Sin AD: posición]{Posición del imán con el modelo no lineal sin atascamiento-deslizamiento (sin AD) y con él (con AD). Abajo, detalle del transitorio entre 9,8 y 12~s.}
    \label{fig:sinad_pos}
\end{figure}

\begin{figure}[ht]
    \centering
    % Estilo común de las figuras sin/con atascamiento-deslizamiento (pgfplots).
\pgfplotsset{
  sinad/.style={
    width=0.92\textwidth, height=5.2cm, xmin=0, xmax=40,
    xlabel={tiempo [\si{s}]},
    grid=major, grid style={gray!25},
    tick label style={font=\small}, label style={font=\small}, legend style={font=\small},
    /pgf/number format/use comma,
    scaled ticks=false,
    y tick label style={/pgf/number format/fixed},
    table/col sep=comma,
    every axis plot/.append style={line join=round},
  },
  detalle/.style={height=4cm, axis background/.style={fill=ctl2!4}},
  sinAD/.style={tray, line width=1.1pt},
  conAD/.style={ctl2, line width=0.7pt},
}
\definecolor{tray}{RGB}{31,94,166}
\definecolor{refc}{RGB}{40,40,40}
\definecolor{ctl2}{RGB}{196,96,40}

\begin{tikzpicture}
\begin{axis}[sinad, ylabel={velocidad [\si{cm/s}]},
             legend style={at={(0.99,0.03)}, anchor=south east}, legend cell align=left]
  \addplot[conAD] table[x=t, y=v_ad] {images/sin_ad_results/tikz/comparacion.csv};
  \addlegendentry{con AD}
  \addplot[sinAD] table[x=t, y=v] {images/sin_ad_results/tikz/comparacion.csv};
  \addlegendentry{sin AD}
\end{axis}
\end{tikzpicture}

    \caption[Sin AD: velocidad]{Velocidad del imán sin y con atascamiento-deslizamiento.}
    \label{fig:sinad_vel}
\end{figure}

\begin{figure}[ht]
    \centering
    % Estilo común de las figuras sin/con atascamiento-deslizamiento (pgfplots).
\pgfplotsset{
  sinad/.style={
    width=0.92\textwidth, height=5.2cm, xmin=0, xmax=40,
    xlabel={tiempo [\si{s}]},
    grid=major, grid style={gray!25},
    tick label style={font=\small}, label style={font=\small}, legend style={font=\small},
    /pgf/number format/use comma,
    scaled ticks=false,
    y tick label style={/pgf/number format/fixed},
    table/col sep=comma,
    every axis plot/.append style={line join=round},
  },
  detalle/.style={height=4cm, axis background/.style={fill=ctl2!4}},
  sinAD/.style={tray, line width=1.1pt},
  conAD/.style={ctl2, line width=0.7pt},
}
\definecolor{tray}{RGB}{31,94,166}
\definecolor{refc}{RGB}{40,40,40}
\definecolor{ctl2}{RGB}{196,96,40}

\begin{tikzpicture}
\begin{axis}[sinad, name=principal, xlabel={}, ylabel={error [\si{cm}]},
             legend style={at={(0.99,0.03)}, anchor=south east}, legend cell align=left]
  \addplot[conAD] table[x=t, y=e_ad] {images/sin_ad_results/tikz/comparacion.csv};
  \addlegendentry{con AD}
  \addplot[sinAD] table[x=t, y=e] {images/sin_ad_results/tikz/comparacion.csv};
  \addlegendentry{sin AD}
  \fill[ctl2, opacity=0.15] ({axis cs:15,0} |- {rel axis cs:0,0}) rectangle ({axis cs:40,0} |- {rel axis cs:0,1});
\end{axis}
\begin{axis}[sinad, detalle, at={(principal.south west)}, anchor=north west, yshift=-0.45cm,
             xmin=15, xmax=40, ymin=-0.03, ymax=0.03, ylabel={detalle [\si{cm}]}]
  \addplot[refc!60, thin] coordinates {(15,0) (40,0)};
  \addplot[conAD] table[x=t, y=e_ad] {images/sin_ad_results/tikz/comparacion.csv};
  \addplot[sinAD] table[x=t, y=e] {images/sin_ad_results/tikz/comparacion.csv};
\end{axis}
\end{tikzpicture}

    \caption[Sin AD: error]{Error de posición sin y con atascamiento-deslizamiento. Abajo, detalle entre 15 y 40~s.}
    \label{fig:sinad_error}
\end{figure}

\begin{figure}[ht]
    \centering
    % Estilo común de las figuras sin/con atascamiento-deslizamiento (pgfplots).
\pgfplotsset{
  sinad/.style={
    width=0.92\textwidth, height=5.2cm, xmin=0, xmax=40,
    xlabel={tiempo [\si{s}]},
    grid=major, grid style={gray!25},
    tick label style={font=\small}, label style={font=\small}, legend style={font=\small},
    /pgf/number format/use comma,
    scaled ticks=false,
    y tick label style={/pgf/number format/fixed},
    table/col sep=comma,
    every axis plot/.append style={line join=round},
  },
  detalle/.style={height=4cm, axis background/.style={fill=ctl2!4}},
  sinAD/.style={tray, line width=1.1pt},
  conAD/.style={ctl2, line width=0.7pt},
}
\definecolor{tray}{RGB}{31,94,166}
\definecolor{refc}{RGB}{40,40,40}
\definecolor{ctl2}{RGB}{196,96,40}

\begin{tikzpicture}
\begin{axis}[sinad, ylabel={señal de control [\si{V}]}, ymin=-0.2, ymax=3.8,
             legend style={at={(0.99,0.03)}, anchor=south east}, legend cell align=left, legend columns=3]
  \addplot[conAD] table[x=t, y=u_ad] {images/sin_ad_results/tikz/comparacion.csv};
  \addlegendentry{con AD}
  \addplot[sinAD] table[x=t, y=u] {images/sin_ad_results/tikz/comparacion.csv};
  \addlegendentry{sin AD}
  \addplot[red!70!black, densely dotted, line width=0.8pt] coordinates {(0,3.5) (40,3.5)};
  \addlegendentry{$u_{\max}$}
\end{axis}
\end{tikzpicture}

    \caption[Sin AD: señal de control]{Señal de control sin y con atascamiento-deslizamiento.}
    \label{fig:sinad_ctrl}
\end{figure}

La figura~\ref{fig:sinad_pos} muestra que, sin AD, el lazo cerrado no lineal alcanza la nueva referencia con un sobrepaso de $23{,}7$\,\% a los $0{,}15$~s del escalón y entra en la banda de $\pm 2$\,\% del cambio de referencia en $0{,}48$~s. Este sobrepaso se encuentra entre los que predice el modelo linealizado en $x_1^*=-3$~cm ($19{,}2$\,\%) y en $x_1^*=-2$~cm ($24{,}7$\,\%), lo que confirma la coherencia entre el modelo empleado para el diseño y la planta no lineal. La velocidad (figura~\ref{fig:sinad_vel}) se anula una vez concluido el transitorio, el error (figura~\ref{fig:sinad_error}) converge a cero y la señal de control (figura~\ref{fig:sinad_ctrl}) converge al voltaje de equilibrio $u_e=2{,}3934$~V correspondiente a $x_1=-3$~cm, sin alcanzar la saturación superior.

La comparación con el caso con AD muestra que el transitorio es prácticamente idéntico (sobrepaso de $23{,}6$\,\%). La diferencia aparece en estado estacionario, donde la fricción estática y la zona muerta sustituyen la convergencia asintótica por un ciclo límite de $0{,}033$~cm pico a pico. Con ello se cumple el tercer objetivo específico de la tesis (capítulo~\ref{chap:introduccion}): simular el desempeño del controlador en el modelo no lineal sin atascamiento-deslizamiento.

\FloatBarrier

\section{Regulación con atascamiento-deslizamiento}
\label{sec:regulacion}

La prueba de regulación aplica el mismo escalón de $-2$ a $-3$~cm en $t=10$~s al lazo completo: PII, zona muerta y fricción de Karnopp, con la señal de control acotada a $u\in[0;\,3{,}5]$~V. Se simulan 40~s con un periodo de muestreo de 1~ms; el controlador parte del estado estacionario correspondiente a $x_1=-2$~cm, de modo que antes del escalón el imán permanece en reposo.

La figura~\ref{posicion} muestra la respuesta. El imán alcanza un mínimo de $-3{,}237$~cm a los $0{,}15$~s del escalón, lo que equivale a un sobrepaso de $23{,}6$\,\%, y entra en la banda de $\pm0{,}05$~cm alrededor de la referencia a los $0{,}40$~s. A partir de ahí no se detiene sobre la referencia: oscila en un ciclo límite (\emph{hunting}, sección~\ref{sec:fund_friccion}) de $0{,}033$~cm pico a pico, con un error máximo de $0{,}018$~cm y un error medio de $6\times10^{-5}$~cm en los últimos 10~s. La velocidad (figura~\ref{velocidad}) alterna pulsos de deslizamiento con intervalos en cero, y el imán permanece atascado el 74\,\% del tiempo en estado estacionario. El error (figura~\ref{error}) refleja el mismo comportamiento.

La señal de control (figura~\ref{senal_control}) cae momentáneamente a 0~V en el instante del escalón, el límite inferior del actuador, para dejar que el imán se aleje de la bobina. Alcanza un máximo de $3{,}28$~V durante el transitorio y en estado estacionario oscila en diente de sierra alrededor de su valor medio de $2{,}41$~V, sin llegar al límite superior. El capítulo~\ref{chap:doble_integral} analiza el origen de esta oscilación.

\begin{figure}[ht]
    \centering
    % Estilo común de las figuras de regulación (pgfplots). Se carga con \input antes de cada figura.
\pgfplotsset{
  regulacion/.style={
    width=0.92\textwidth, height=5.6cm,
    xmin=0, xmax=40,
    xlabel={tiempo [\si{s}]},
    grid=major, grid style={gray!25},
    tick label style={font=\small}, label style={font=\small}, legend style={font=\small},
    /pgf/number format/use comma,
    scaled y ticks=false,
    y tick label style={/pgf/number format/fixed},
    table/col sep=comma,
    every axis plot/.append style={line join=round},
  },
}
\definecolor{tray}{RGB}{31,94,166}
\definecolor{refc}{RGB}{40,40,40}
\definecolor{ctl2}{RGB}{196,96,40}

\begin{tikzpicture}
\begin{axis}[regulacion, ylabel={posición [\si{cm}]}, ymin=-3.35, ymax=-1.85,
             legend pos=south west, legend cell align=left]
  \addplot[tray, line width=1pt] table[x=t, y=x] {images/regulation_results/tikz/posicion.csv};
  \addlegendentry{posición $x_1$}
  \addplot[refc, densely dashed, line width=0.8pt] table[x=t, y=r] {images/regulation_results/tikz/posicion.csv};
  \addlegendentry{referencia}
  \coordinate (zoomA) at (axis cs:9.8,-3.27);
  \coordinate (zoomB) at (axis cs:14,-2.95);
  \coordinate (inset) at (axis cs:39.2,-1.93);
\end{axis}
\draw[gray, thin] (zoomA) rectangle (zoomB);
\begin{axis}[regulacion, at={(inset)}, anchor=north east, width=0.5\textwidth, height=3.1cm,
             xmin=9.8, xmax=14, ymin=-3.27, ymax=-2.95, xlabel={}, ylabel={},
             tick label style={font=\scriptsize}, axis background/.style={fill=white},
             xtick={10,11,12,13,14}, ytick={-3.2,-3.1,-3}, name=zoom]
  \addplot[tray, line width=0.9pt] table[x=t, y=x] {images/regulation_results/tikz/posicion.csv};
  \addplot[refc, densely dashed, line width=0.7pt] table[x=t, y=r] {images/regulation_results/tikz/posicion.csv};
\end{axis}
\draw[gray, thin, densely dotted] (zoomB) -- (zoom.south west);
\end{tikzpicture}

    \caption[Posición (regulación)]{Posición del imán ante un escalón de referencia de $-2$ a $-3$~cm en $t=10$~s (PII, $u\in[0;\,3{,}5]$~V). El recuadro amplía el transitorio y el ciclo límite de atascamiento-deslizamiento.}
    \label{posicion}
\end{figure}

\begin{figure}[ht]
    \centering
    % Estilo común de las figuras de regulación (pgfplots). Se carga con \input antes de cada figura.
\pgfplotsset{
  regulacion/.style={
    width=0.92\textwidth, height=5.6cm,
    xmin=0, xmax=40,
    xlabel={tiempo [\si{s}]},
    grid=major, grid style={gray!25},
    tick label style={font=\small}, label style={font=\small}, legend style={font=\small},
    /pgf/number format/use comma,
    scaled y ticks=false,
    y tick label style={/pgf/number format/fixed},
    table/col sep=comma,
    every axis plot/.append style={line join=round},
  },
}
\definecolor{tray}{RGB}{31,94,166}
\definecolor{refc}{RGB}{40,40,40}
\definecolor{ctl2}{RGB}{196,96,40}

\begin{tikzpicture}
\begin{axis}[regulacion, ylabel={velocidad [\si{cm/s}]}]
  \addplot[tray, line width=0.8pt] table[x=t, y=v] {images/regulation_results/tikz/velocidad.csv};
\end{axis}
\end{tikzpicture}

    \caption[Velocidad (regulación)]{Velocidad del imán. Los tramos en cero corresponden al régimen de atascamiento.}
    \label{velocidad}
\end{figure}

\begin{figure}[ht]
    \centering
    % Estilo común de las figuras de regulación (pgfplots). Se carga con \input antes de cada figura.
\pgfplotsset{
  regulacion/.style={
    width=0.92\textwidth, height=5.6cm,
    xmin=0, xmax=40,
    xlabel={tiempo [\si{s}]},
    grid=major, grid style={gray!25},
    tick label style={font=\small}, label style={font=\small}, legend style={font=\small},
    /pgf/number format/use comma,
    scaled y ticks=false,
    y tick label style={/pgf/number format/fixed},
    table/col sep=comma,
    every axis plot/.append style={line join=round},
  },
}
\definecolor{tray}{RGB}{31,94,166}
\definecolor{refc}{RGB}{40,40,40}
\definecolor{ctl2}{RGB}{196,96,40}

\begin{tikzpicture}
\begin{axis}[regulacion, ylabel={error $e=r-x_1$ [\si{cm}]}]
  \addplot[refc!60, thin] coordinates {(0,0) (40,0)};
  \addplot[tray, line width=0.9pt] table[x=t, y=e] {images/regulation_results/tikz/error.csv};
\end{axis}
\end{tikzpicture}

    \caption[Error (regulación)]{Error de posición $e=r-x_1$.}
    \label{error}
\end{figure}

\begin{figure}[ht]
    \centering
    % Estilo común de las figuras de regulación (pgfplots). Se carga con \input antes de cada figura.
\pgfplotsset{
  regulacion/.style={
    width=0.92\textwidth, height=5.6cm,
    xmin=0, xmax=40,
    xlabel={tiempo [\si{s}]},
    grid=major, grid style={gray!25},
    tick label style={font=\small}, label style={font=\small}, legend style={font=\small},
    /pgf/number format/use comma,
    scaled y ticks=false,
    y tick label style={/pgf/number format/fixed},
    table/col sep=comma,
    every axis plot/.append style={line join=round},
  },
}
\definecolor{tray}{RGB}{31,94,166}
\definecolor{refc}{RGB}{40,40,40}
\definecolor{ctl2}{RGB}{196,96,40}

\begin{tikzpicture}
\begin{axis}[regulacion, ylabel={señal de control [\si{V}]}, ymin=-0.2, ymax=3.8,
             legend pos=south east, legend cell align=left]
  \addplot[ctl2!45, line width=1.6pt] table[x=t, y=u_antes] {images/regulation_results/tikz/control.csv};
  \addlegendentry{antes de la zona muerta}
  \addplot[tray, line width=0.7pt] table[x=t, y=u_desp] {images/regulation_results/tikz/control.csv};
  \addlegendentry{después de la zona muerta}
  \addplot[red!70!black, densely dotted, line width=0.8pt] coordinates {(0,3.5) (40,3.5)};
  \addlegendentry{$u_{\max}=\SI{3,5}{V}$}
\end{axis}
\end{tikzpicture}

    \caption[Señal de control (regulación)]{Señal de control antes y después de la zona muerta.}
    \label{senal_control}
\end{figure}

\FloatBarrier

\section{Causa del ciclo límite: modelo lineal con atascamiento-deslizamiento}
\label{sec:lineal_ad}
% Generado a partir de calculus/Final_Bien/evidencia_P8.m (simulador corregido).

El ciclo límite podría deberse a la no linealidad de la fuerza electromagnética o a las no linealidades no suaves del actuador y la fricción. Para separarlas, el PII se probó sobre el modelo linealizado en $x_1^*=-2{,}5$~cm al que se añadieron la zona muerta y la fricción de Karnopp, con la misma prueba de la sección~\ref{sec:regulacion} ($-2\to-3$~cm, $u\in[0;\,3{,}5]$~V). La figura~\ref{fig:p8_lineal_ad} compara ambos resultados. El comportamiento es prácticamente el mismo: sobrepaso de $24{,}4$\,\% frente a $23{,}6$\,\%, entrada en la banda de $\pm5$\,\% en $0{,}41$ frente a $0{,}40$~s y un ciclo límite de $0{,}028$ frente a $0{,}033$~cm pico a pico. La única diferencia apreciable es el voltaje de equilibrio final ($2{,}36$~V en el modelo lineal frente a $2{,}39$~V), consecuencia del error de linealización a $0{,}5$~cm del punto de operación.

El ciclo límite aparece, pues, aun cuando la fuerza electromagnética es lineal. Se debe a la zona muerta y a la fricción estática, no a la no linealidad de la fuerza.

\begin{figure}[ht]
    \centering
    % Estilo común de las figuras de evidencia P8 (pgfplots).
\pgfplotsset{
  pocho/.style={
    width=0.92\textwidth, height=5.2cm,
    grid=major, grid style={gray!25},
    tick label style={font=\small}, label style={font=\small}, legend style={font=\small},
    /pgf/number format/use comma,
    scaled ticks=false,
    y tick label style={/pgf/number format/fixed},
    table/col sep=comma, unbounded coords=jump,
    every axis plot/.append style={line join=round},
  },
}
\definecolor{tray}{RGB}{31,94,166}
\definecolor{refc}{RGB}{40,40,40}
\definecolor{ctl2}{RGB}{196,96,40}
\definecolor{ver}{RGB}{46,133,64}

\begin{tikzpicture}
\begin{axis}[pocho, name=principal, xmin=0, xmax=40, ymin=-3.35, ymax=-1.85, xlabel={},
    ylabel={posición [\si{cm}]}, legend style={at={(0.99,0.97)}, anchor=north east}, legend cell align=left]
  \addplot[ctl2, line width=0.7pt] table[x=t, y=x_nl] {images/p8_results/tikz/lineal_ad.csv}; \addlegendentry{no lineal con AD}
  \addplot[tray, line width=1pt] table[x=t, y=x_lin] {images/p8_results/tikz/lineal_ad.csv}; \addlegendentry{lineal con AD}
\end{axis}
\begin{axis}[pocho, at={(principal.south west)}, anchor=north west, yshift=-0.45cm, height=4.4cm,
    xmin=0, xmax=40, ymin=-0.2, ymax=3.8, xlabel={tiempo [\si{s}]}, ylabel={$u$ [\si{V}]}]
  \addplot[ctl2, line width=0.7pt] table[x=t, y=u_nl] {images/p8_results/tikz/lineal_ad.csv};
  \addplot[tray, line width=0.7pt] table[x=t, y=u_lin] {images/p8_results/tikz/lineal_ad.csv};
\end{axis}
\end{tikzpicture}

    \caption[Modelo lineal con atascamiento-deslizamiento]{Regulación con zona muerta y fricción de Karnopp sobre el modelo lineal (en $x_1^*=-2{,}5$~cm) y sobre el modelo no lineal. Arriba, posición; abajo, señal de control.}
    \label{fig:p8_lineal_ad}
\end{figure}

\FloatBarrier

\section{Seguimiento}

La regulación fija la referencia durante decenas de segundos. Las pruebas de seguimiento, en cambio, hacen que la referencia varíe continuamente, de modo que el imán debe romper la fricción estática en ambos sentidos a lo largo de toda la prueba. Se mantiene el intervalo de $-3$ a $-2$~cm, en el que el actuador limitado a $u\in[0;\,3{,}5]$~V puede vencer la fricción estática en ambos sentidos. El antecedente más cercano, el lazo externo con acción integral de Hernández Alcántara et al.~\cite{alcantara} en la configuración repulsiva (estable) del mismo dispositivo, se revisó en la sección~\ref{sec:motivacion}; aquí la tarea es más exigente porque, en configuración atractiva, el controlador debe además estabilizar la planta.

\subsection{Seguimiento de referencia senoidal}

La referencia de posición tiene la forma
\begin{equation}
  r(t) = x_0 + A\,\sin(2\pi f\, t) ,
  \label{eq:ref_senoidal}
\end{equation}
con valor nominal $x_0 = -2{,}5$~cm, amplitud $A = 0{,}5$~cm y frecuencia $f = 0{,}001$~Hz (periodo $T = 1000$~s). La amplitud hace que la referencia excursione simétricamente entre $-3$ y $-2$~cm, el mismo intervalo de la prueba de regulación. Se simularon dos periodos (2000~s).

Sin atascamiento-deslizamiento, el lazo cerrado lineal con función de transferencia $T(s)$ y $T(j\omega_0)=|T|\,e^{j\varphi}$ produciría, ante una referencia de frecuencia $\omega_0 = 2\pi f$, el error en estado estacionario
\begin{equation}
  e_{ss}(t) = A\bigl[1 - |T|\cos\varphi\bigr]\sin(\omega_0 t)
              - A\,|T|\sin\varphi\,\cos(\omega_0 t) ,
  \label{eq:ess_senoidal}
\end{equation}
que es pequeño cuando $T(j\omega_0)$ es próximo a la unidad, como ocurre con una frecuencia tan baja. Con atascamiento-deslizamiento, sin embargo, el error real lo domina el ciclo límite y no esta componente lineal.

\begin{figure}[ht]
    \centering
    % Estilo común de las figuras de seguimiento (pgfplots).
\pgfplotsset{
  seguimiento/.style={
    width=0.92\textwidth, height=5.4cm,
    xlabel={tiempo [\si{s}]},
    grid=major, grid style={gray!25},
    tick label style={font=\small}, label style={font=\small}, legend style={font=\small},
    /pgf/number format/use comma,
    scaled ticks=false,
    y tick label style={/pgf/number format/fixed},
    x tick label style={/pgf/number format/fixed, /pgf/number format/1000 sep={}},
    table/col sep=comma,
    every axis plot/.append style={line join=round},
  },
}
\definecolor{tray}{RGB}{31,94,166}
\definecolor{refc}{RGB}{40,40,40}
\definecolor{ctl2}{RGB}{196,96,40}

\begin{tikzpicture}
\begin{axis}[seguimiento, name=principal, xmin=0, xmax=2000, ymin=-3.15, ymax=-1.85,
             xlabel={}, ylabel={posición [\si{cm}]},
             legend style={at={(0.99,0.03)}, anchor=south east}, legend cell align=left]
  \addplot[tray, line width=1pt] table[x=t, y=x] {images/seguimiento_results/tikz/sen_posicion.csv};
  \addlegendentry{posición $x_1$}
  \addplot[refc, densely dashed, line width=0.8pt] table[x=t, y=r] {images/seguimiento_results/tikz/sen_posicion.csv};
  \addlegendentry{referencia $r$}
  \fill[ctl2, opacity=0.15] (axis cs:244,-3.15) rectangle (axis cs:256,-1.85);
\end{axis}
\begin{axis}[seguimiento, at={(principal.south west)}, anchor=north west, yshift=-0.45cm,
             height=4.2cm, xmin=244, xmax=256,
             ylabel={detalle [\si{cm}]}, axis background/.style={fill=ctl2!4}]
  \addplot[tray, line width=1pt] table[x=t, y=x] {images/seguimiento_results/tikz/sen_detalle.csv};
  \addplot[refc, densely dashed, line width=0.8pt] table[x=t, y=r] {images/seguimiento_results/tikz/sen_detalle.csv};
\end{axis}
\end{tikzpicture}

    \caption[Seguimiento senoidal: posición]{Seguimiento de la referencia senoidal. Arriba, posición del imán (línea continua) y referencia (discontinua) durante dos periodos; la banda sombreada marca la ventana de detalle de abajo (244--256~s, máximo de la referencia), donde se observa el ciclo límite de atascamiento-deslizamiento alrededor de la referencia.}
    \label{fig:sen_pos}
\end{figure}

La figura~\ref{fig:sen_pos} muestra que, a la escala de la referencia, la posición sigue a la senoide sin desfase ni atenuación apreciables, en concordancia con $|T(j\omega_0)|\approx 1$. La ventana de detalle revela, sin embargo, que el imán avanza en saltos y permanece atascado entre ellos, formando un ciclo límite de aproximadamente $0{,}02$~cm pico a pico y periodo cercano a 1~s.

\FloatBarrier

\begin{figure}[ht]
    \centering
    % Estilo común de las figuras de seguimiento (pgfplots).
\pgfplotsset{
  seguimiento/.style={
    width=0.92\textwidth, height=5.4cm,
    xlabel={tiempo [\si{s}]},
    grid=major, grid style={gray!25},
    tick label style={font=\small}, label style={font=\small}, legend style={font=\small},
    /pgf/number format/use comma,
    scaled ticks=false,
    y tick label style={/pgf/number format/fixed},
    x tick label style={/pgf/number format/fixed, /pgf/number format/1000 sep={}},
    table/col sep=comma,
    every axis plot/.append style={line join=round},
  },
}
\definecolor{tray}{RGB}{31,94,166}
\definecolor{refc}{RGB}{40,40,40}
\definecolor{ctl2}{RGB}{196,96,40}

\begin{tikzpicture}
\begin{axis}[seguimiento, name=principal, xmin=0, xmax=2000, xlabel={}, ylabel={velocidad [\si{cm/s}]}]
  \addplot[tray, line width=0.6pt] table[x=t, y=v] {images/seguimiento_results/tikz/sen_velocidad.csv};
  \fill[ctl2, opacity=0.15] ({axis cs:244,0} |- {rel axis cs:0,0}) rectangle ({axis cs:256,0} |- {rel axis cs:0,1});
\end{axis}
\begin{axis}[seguimiento, at={(principal.south west)}, anchor=north west, yshift=-0.45cm,
             height=4.2cm, xmin=244, xmax=256, ylabel={detalle [\si{cm/s}]}, axis background/.style={fill=ctl2!4}]
  \addplot[tray, line width=0.8pt] table[x=t, y=v] {images/seguimiento_results/tikz/sen_detalle.csv};
\end{axis}
\end{tikzpicture}

    \caption[Seguimiento senoidal: velocidad]{Velocidad del imán durante el seguimiento senoidal. Los pulsos corresponden a los eventos de deslizamiento y los tramos en cero al régimen de atascamiento (75\,\% del tiempo).}
    \label{fig:sen_vel}
\end{figure}

La figura~\ref{fig:sen_vel} muestra que la velocidad alterna pulsos de deslizamiento de hasta $\pm 1$~cm/s con tramos de velocidad nula. En los 2000~s simulados se registraron 4072 rupturas (\emph{breakaway}; unas dos por segundo) y el imán permaneció atascado el 75\,\% del tiempo. Las transiciones entre atascamiento y deslizamiento ocurren de forma continua, también lejos de los instantes en que la referencia invierte su sentido ($t = T/4,\,3T/4,\ldots$). La condición que el modelo de Karnopp aplica para declarar el atascamiento se presentó en la sección~\ref{sec:friccion_modelo}.

\FloatBarrier

\begin{figure}[ht]
    \centering
    % Estilo común de las figuras de seguimiento (pgfplots).
\pgfplotsset{
  seguimiento/.style={
    width=0.92\textwidth, height=5.4cm,
    xlabel={tiempo [\si{s}]},
    grid=major, grid style={gray!25},
    tick label style={font=\small}, label style={font=\small}, legend style={font=\small},
    /pgf/number format/use comma,
    scaled ticks=false,
    y tick label style={/pgf/number format/fixed},
    x tick label style={/pgf/number format/fixed, /pgf/number format/1000 sep={}},
    table/col sep=comma,
    every axis plot/.append style={line join=round},
  },
}
\definecolor{tray}{RGB}{31,94,166}
\definecolor{refc}{RGB}{40,40,40}
\definecolor{ctl2}{RGB}{196,96,40}

\begin{tikzpicture}
\begin{axis}[seguimiento, xmin=0, xmax=2000, ylabel={error $e=r-x_1$ [\si{cm}]}]
  \addplot[refc!60, thin] coordinates {(0,0) (2000,0)};
  \addplot[tray, line width=0.7pt] table[x=t, y=e] {images/seguimiento_results/tikz/sen_error.csv};
\end{axis}
\end{tikzpicture}

    \caption[Seguimiento senoidal: error]{Error de seguimiento $e(t)=r(t)-x_1(t)$ para la referencia senoidal.}
    \label{fig:sen_error}
\end{figure}

La figura~\ref{fig:sen_error} muestra que el error permanece acotado en una banda de $\pm 0{,}02$~cm (error máximo $0{,}0198$~cm, valor eficaz $0{,}0124$~cm) con media prácticamente nula. La amplitud de la banda varía ligeramente con la posición de la referencia, pero no se observan picos aislados: el error está dominado por el ciclo límite y no por la dinámica lineal de seguimiento.

\FloatBarrier

\begin{figure}[ht]
    \centering
    % Estilo común de las figuras de seguimiento (pgfplots).
\pgfplotsset{
  seguimiento/.style={
    width=0.92\textwidth, height=5.4cm,
    xlabel={tiempo [\si{s}]},
    grid=major, grid style={gray!25},
    tick label style={font=\small}, label style={font=\small}, legend style={font=\small},
    /pgf/number format/use comma,
    scaled ticks=false,
    y tick label style={/pgf/number format/fixed},
    x tick label style={/pgf/number format/fixed, /pgf/number format/1000 sep={}},
    table/col sep=comma,
    every axis plot/.append style={line join=round},
  },
}
\definecolor{tray}{RGB}{31,94,166}
\definecolor{refc}{RGB}{40,40,40}
\definecolor{ctl2}{RGB}{196,96,40}

\begin{tikzpicture}
\begin{axis}[seguimiento, name=principal, xmin=0, xmax=2000, ymin=-0.2, ymax=3.8, xlabel={},
             ylabel={señal de control [\si{V}]},
             legend style={at={(0.99,0.03)}, anchor=south east}, legend cell align=left, legend columns=3]
  \addplot[ctl2!45, line width=1.2pt] table[x=t, y=u_antes] {images/seguimiento_results/tikz/sen_control.csv};
  \addlegendentry{antes de la ZM}
  \addplot[tray, line width=0.5pt] table[x=t, y=u_desp] {images/seguimiento_results/tikz/sen_control.csv};
  \addlegendentry{después de la ZM}
  \addplot[red!70!black, densely dotted, line width=0.8pt] coordinates {(0,3.5) (2000,3.5)};
  \addlegendentry{$u_{\max}$}
  \fill[ctl2, opacity=0.15] (axis cs:244,-0.2) rectangle (axis cs:256,3.8);
\end{axis}
\begin{axis}[seguimiento, at={(principal.south west)}, anchor=north west, yshift=-0.45cm,
             height=4.2cm, xmin=244, xmax=256, ylabel={detalle [\si{V}]}, axis background/.style={fill=ctl2!4}]
  \addplot[tray, line width=0.8pt] table[x=t, y=u] {images/seguimiento_results/tikz/sen_detalle.csv};
\end{axis}
\end{tikzpicture}

    \caption[Seguimiento senoidal: señal de control]{Voltaje de control antes y después de la zona muerta durante el seguimiento senoidal. Abajo, detalle de 244 a 256~s: perfil en diente de sierra del voltaje en estado estacionario.}
    \label{fig:sen_ctrl}
\end{figure}

La figura~\ref{fig:sen_ctrl} muestra que el voltaje sigue la forma de la referencia entre $1{,}33$ y $2{,}78$~V, dentro del rango admisible. El detalle revela un perfil en diente de sierra: el voltaje crece durante cada atascamiento hasta que el imán se libera y cae bruscamente tras la ruptura.

\FloatBarrier

\subsection{Seguimiento de referencia trapezoidal}

La referencia trapezoidal alterna segmentos planos en $x_1 = -2$~cm y $x_1 = -3$~cm, conectados por rampas lineales de $250$~s de duración cada una, con velocidad constante $\dot{r} = (x_{f}-x_{i})/250$~cm/s. Combina así fases de regulación (los segmentos planos) con fases de transición (las rampas), y es más exigente que la senoidal porque su derivada es discontinua al inicio y al fin de cada rampa.

Según el modelo lineal, durante la rampa el error de seguimiento converge a un valor constante que fija el tipo del lazo (sección~\ref{sec:fund_tipo}). La segunda acción integral elimina el error de posición constante que un controlador de tipo uno mantendría; para la pendiente empleada, ese error sería del orden de $10^{-5}$~cm, muy inferior al ciclo límite, y el capítulo~\ref{chap:doble_integral} lo compara con el del PI. En los segmentos planos, sin atascamiento-deslizamiento el error se anularía; con él, como en la regulación, el error medio se anula pero el imán oscila alrededor de la referencia.

Al inicio de cada rampa el imán debe abandonar el atascamiento en que quedó durante el segmento plano previo: la fuerza que el controlador acumula tiene que superar la fricción estática máxima $F_s = 20$~kg\,cm/s$^2$ antes de que el imán se desplace. En la simulación, este retraso es de aproximadamente un segundo al inicio de la primera rampa y no produce un pico de error mayor que la banda del ciclo límite.

\begin{figure}[ht]
    \centering
    % Estilo común de las figuras de seguimiento (pgfplots).
\pgfplotsset{
  seguimiento/.style={
    width=0.92\textwidth, height=5.4cm,
    xlabel={tiempo [\si{s}]},
    grid=major, grid style={gray!25},
    tick label style={font=\small}, label style={font=\small}, legend style={font=\small},
    /pgf/number format/use comma,
    scaled ticks=false,
    y tick label style={/pgf/number format/fixed},
    x tick label style={/pgf/number format/fixed, /pgf/number format/1000 sep={}},
    table/col sep=comma,
    every axis plot/.append style={line join=round},
  },
}
\definecolor{tray}{RGB}{31,94,166}
\definecolor{refc}{RGB}{40,40,40}
\definecolor{ctl2}{RGB}{196,96,40}

\begin{tikzpicture}
\begin{axis}[seguimiento, name=principal, xmin=0, xmax=2500, ymin=-3.15, ymax=-1.85,
             xlabel={}, ylabel={posición [\si{cm}]},
             legend style={at={(0.99,0.03)}, anchor=south east}, legend cell align=left]
  \addplot[tray, line width=1pt] table[x=t, y=x] {images/seguimiento_results/tikz/trap_posicion.csv};
  \addlegendentry{posición $x_1$}
  \addplot[refc, densely dashed, line width=0.8pt] table[x=t, y=r] {images/seguimiento_results/tikz/trap_posicion.csv};
  \addlegendentry{referencia $r$}
  \fill[ctl2, opacity=0.15] (axis cs:246,-3.15) rectangle (axis cs:266,-1.85);
\end{axis}
\begin{axis}[seguimiento, at={(principal.south west)}, anchor=north west, yshift=-0.45cm,
             height=4.2cm, xmin=246, xmax=266,
             ylabel={detalle [\si{cm}]}, axis background/.style={fill=ctl2!4}]
  \addplot[tray, line width=1pt] table[x=t, y=x] {images/seguimiento_results/tikz/trap_detalle.csv};
  \addplot[refc, densely dashed, line width=0.8pt] table[x=t, y=r] {images/seguimiento_results/tikz/trap_detalle.csv};
\end{axis}
\end{tikzpicture}

    \caption[Seguimiento trapezoidal: posición]{Seguimiento de la referencia trapezoidal. Arriba, posición y referencia durante dos periodos; abajo, detalle de 246 a 266~s, que incluye el inicio de la primera rampa en $t=250$~s.}
    \label{fig:trap_pos}
\end{figure}

La figura~\ref{fig:trap_pos} muestra que el imán sigue las rampas y los segmentos planos sin error apreciable a la escala de la referencia. El detalle del inicio de la primera rampa muestra que el imán permanece atascado alrededor de un segundo después de que la referencia comienza a moverse, y a partir de ahí sigue la rampa en escalones sucesivos de atascamiento-deslizamiento.

\FloatBarrier

\begin{figure}[ht]
    \centering
    % Estilo común de las figuras de seguimiento (pgfplots).
\pgfplotsset{
  seguimiento/.style={
    width=0.92\textwidth, height=5.4cm,
    xlabel={tiempo [\si{s}]},
    grid=major, grid style={gray!25},
    tick label style={font=\small}, label style={font=\small}, legend style={font=\small},
    /pgf/number format/use comma,
    scaled ticks=false,
    y tick label style={/pgf/number format/fixed},
    x tick label style={/pgf/number format/fixed, /pgf/number format/1000 sep={}},
    table/col sep=comma,
    every axis plot/.append style={line join=round},
  },
}
\definecolor{tray}{RGB}{31,94,166}
\definecolor{refc}{RGB}{40,40,40}
\definecolor{ctl2}{RGB}{196,96,40}

\begin{tikzpicture}
\begin{axis}[seguimiento, name=principal, xmin=0, xmax=2500, xlabel={}, ylabel={velocidad [\si{cm/s}]}]
  \addplot[tray, line width=0.6pt] table[x=t, y=v] {images/seguimiento_results/tikz/trap_velocidad.csv};
  \fill[ctl2, opacity=0.15] ({axis cs:246,0} |- {rel axis cs:0,0}) rectangle ({axis cs:266,0} |- {rel axis cs:0,1});
\end{axis}
\begin{axis}[seguimiento, at={(principal.south west)}, anchor=north west, yshift=-0.45cm,
             height=4.2cm, xmin=246, xmax=266, ylabel={detalle [\si{cm/s}]}, axis background/.style={fill=ctl2!4}]
  \addplot[tray, line width=0.8pt] table[x=t, y=v] {images/seguimiento_results/tikz/trap_detalle.csv};
\end{axis}
\end{tikzpicture}

    \caption[Seguimiento trapezoidal: velocidad]{Velocidad del imán durante el seguimiento trapezoidal. El imán permanece en reposo durante el primer segmento plano; a partir de la primera rampa, los pulsos de deslizamiento se mantienen durante el resto de la prueba.}
    \label{fig:trap_vel}
\end{figure}

La figura~\ref{fig:trap_vel} muestra que el imán permanece en reposo durante el primer segmento plano, en el que parte del equilibrio. A partir de la primera rampa aparecen pulsos de deslizamiento de hasta $\pm 1$~cm/s, que continúan también en los segmentos planos posteriores: una vez iniciado, el ciclo límite no se extingue. En total se registraron 4400 rupturas en los 2500~s (dos periodos) y el imán estuvo atascado el 78\,\% del tiempo.

\FloatBarrier

\begin{figure}[ht]
    \centering
    % Estilo común de las figuras de seguimiento (pgfplots).
\pgfplotsset{
  seguimiento/.style={
    width=0.92\textwidth, height=5.4cm,
    xlabel={tiempo [\si{s}]},
    grid=major, grid style={gray!25},
    tick label style={font=\small}, label style={font=\small}, legend style={font=\small},
    /pgf/number format/use comma,
    scaled ticks=false,
    y tick label style={/pgf/number format/fixed},
    x tick label style={/pgf/number format/fixed, /pgf/number format/1000 sep={}},
    table/col sep=comma,
    every axis plot/.append style={line join=round},
  },
}
\definecolor{tray}{RGB}{31,94,166}
\definecolor{refc}{RGB}{40,40,40}
\definecolor{ctl2}{RGB}{196,96,40}

\begin{tikzpicture}
\begin{axis}[seguimiento, xmin=0, xmax=2500, ylabel={error $e=r-x_1$ [\si{cm}]}]
  \addplot[refc!60, thin] coordinates {(0,0) (2500,0)};
  \addplot[tray, line width=0.7pt] table[x=t, y=e] {images/seguimiento_results/tikz/trap_error.csv};
\end{axis}
\end{tikzpicture}

    \caption[Seguimiento trapezoidal: error]{Error de seguimiento $e(t)=r(t)-x_1(t)$ para la referencia trapezoidal.}
    \label{fig:trap_error}
\end{figure}

La figura~\ref{fig:trap_error} muestra un error acotado en $\pm 0{,}022$~cm (valor eficaz $0{,}0114$~cm) y media nula, tanto en las rampas como en los segmentos planos. Al inicio de las rampas el error conserva la banda de amplitud casi uniforme que impone el ciclo límite, sin picos aislados.

\FloatBarrier

\begin{figure}[ht]
    \centering
    % Estilo común de las figuras de seguimiento (pgfplots).
\pgfplotsset{
  seguimiento/.style={
    width=0.92\textwidth, height=5.4cm,
    xlabel={tiempo [\si{s}]},
    grid=major, grid style={gray!25},
    tick label style={font=\small}, label style={font=\small}, legend style={font=\small},
    /pgf/number format/use comma,
    scaled ticks=false,
    y tick label style={/pgf/number format/fixed},
    x tick label style={/pgf/number format/fixed, /pgf/number format/1000 sep={}},
    table/col sep=comma,
    every axis plot/.append style={line join=round},
  },
}
\definecolor{tray}{RGB}{31,94,166}
\definecolor{refc}{RGB}{40,40,40}
\definecolor{ctl2}{RGB}{196,96,40}

\begin{tikzpicture}
\begin{axis}[seguimiento, name=principal, xmin=0, xmax=2500, ymin=-0.2, ymax=3.8, xlabel={},
             ylabel={señal de control [\si{V}]},
             legend style={at={(0.99,0.03)}, anchor=south east}, legend cell align=left, legend columns=3]
  \addplot[ctl2!45, line width=1.2pt] table[x=t, y=u_antes] {images/seguimiento_results/tikz/trap_control.csv};
  \addlegendentry{antes de la ZM}
  \addplot[tray, line width=0.5pt] table[x=t, y=u_desp] {images/seguimiento_results/tikz/trap_control.csv};
  \addlegendentry{después de la ZM}
  \addplot[red!70!black, densely dotted, line width=0.8pt] coordinates {(0,3.5) (2500,3.5)};
  \addlegendentry{$u_{\max}$}
  \fill[ctl2, opacity=0.15] (axis cs:246,-0.2) rectangle (axis cs:266,3.8);
\end{axis}
\begin{axis}[seguimiento, at={(principal.south west)}, anchor=north west, yshift=-0.45cm,
             height=4.2cm, xmin=246, xmax=266, ylabel={detalle [\si{V}]}, axis background/.style={fill=ctl2!4}]
  \addplot[tray, line width=0.8pt] table[x=t, y=u] {images/seguimiento_results/tikz/trap_detalle.csv};
\end{axis}
\end{tikzpicture}

    \caption[Seguimiento trapezoidal: señal de control]{Voltaje de control antes y después de la zona muerta durante el seguimiento trapezoidal. Abajo, detalle de 246 a 266~s.}
    \label{fig:trap_ctrl}
\end{figure}

La figura~\ref{fig:trap_ctrl} muestra que el voltaje sigue el perfil trapezoidal entre $1{,}32$ y $2{,}78$~V, sin saturar. El detalle muestra cómo, al inicio de la rampa, el voltaje desciende de forma suave hasta que la fuerza neta vence la fricción estática, y a partir de ese instante adopta el mismo perfil en diente de sierra de la prueba senoidal.

\FloatBarrier

% Generado a partir de calculus/Final_Bien/evidencia_P8.m (simulador corregido).
\section{Configuración estable}

Queda por verificar si el controlador diseñado para la configuración inestable sirve también en la estable. En la configuración repulsiva actúa la bobina inferior. Con $x$ medida hacia arriba, la distancia entre el imán y esa bobina es $a+x$, la fuerza electromagnética es $u/[b(a+x)^4]$ y el modelo linealizado tiene la matriz de estado $\bigl[\begin{smallmatrix}0&1\\-4g/(a+x_1^*)&-c_1\end{smallmatrix}\bigr]$, cuyos polos son complejos con parte real negativa: $-4{,}28\pm21{,}49j$ en $x_1^*=1$~cm y $-4{,}28\pm20{,}24j$ en $x_1^*=2$~cm (factor de amortiguamiento cercano a $0{,}2$). Se evaluaron tres controladores diseñados durante el trabajo: el PI de comparación (sección~\ref{sec:pi}), una versión preliminar del PII y el PII definitivo $C(s)$ de la sección~\ref{sec:pii}. Con la planta linealizada en $x_1^*=1{,}5$~cm, los tres estabilizan el lazo, con márgenes de fase de $50{,}7^\circ$, $51{,}0^\circ$ y $43{,}3^\circ$, respectivamente. El mismo controlador diseñado para la configuración inestable estabiliza, por tanto, también la estable, aunque con un margen de fase menor.

\begin{table}[ht]
\centering
\caption{Configuración estable: escalón de 1 a 2~cm con zona muerta, fricción de Karnopp y $u\in[0;\,3{,}5]$~V.}
\label{tab:p8_estable}
\begin{tabular}{|l|c|c|c|}
\hline
 & PI & PII preliminar & PII (tesis) \\ \hline
Margen de fase (lineal, $x_1^*=1{,}5$~cm) & $50{,}7^\circ$ & $51{,}0^\circ$ & $43{,}3^\circ$ \\ \hline
Sobrepaso (lineal) & $19{,}7$\,\% & $18{,}6$\,\% & $26{,}8$\,\% \\ \hline
Sobrepaso (no lineal con AD) & $10{,}5$\,\% & $11{,}1$\,\% & $9{,}8$\,\% \\ \hline
Entrada en la banda de $\pm5$\,\% & $0{,}24$~s & $0{,}36$~s & $0{,}35$~s \\ \hline
Ciclo límite pico a pico & $0{,}018$~cm & $0{,}023$~cm & $0{,}020$~cm \\ \hline
\end{tabular}
\end{table}

La figura~\ref{fig:p8_estable} y la tabla~\ref{tab:p8_estable} muestran que los tres controladores llevan el imán de 1 a 2~cm y que, como en la configuración atractiva, la zona muerta y la fricción estática producen un ciclo límite en estado estacionario, aquí de menor amplitud. En el transitorio la señal de control alcanza brevemente la saturación de $3{,}5$~V.

\begin{figure}[ht]
    \centering
    % Estilo común de las figuras de evidencia P8 (pgfplots).
\pgfplotsset{
  pocho/.style={
    width=0.92\textwidth, height=5.2cm,
    grid=major, grid style={gray!25},
    tick label style={font=\small}, label style={font=\small}, legend style={font=\small},
    /pgf/number format/use comma,
    scaled ticks=false,
    y tick label style={/pgf/number format/fixed},
    table/col sep=comma, unbounded coords=jump,
    every axis plot/.append style={line join=round},
  },
}
\definecolor{tray}{RGB}{31,94,166}
\definecolor{refc}{RGB}{40,40,40}
\definecolor{ctl2}{RGB}{196,96,40}
\definecolor{ver}{RGB}{46,133,64}

\begin{tikzpicture}
\begin{axis}[pocho, name=principal, xmin=0, xmax=40, ymin=0.8, ymax=2.3, xlabel={},
    ylabel={posición [\si{cm}]}, legend style={at={(0.99,0.03)}, anchor=south east}, legend cell align=left]
  \addplot[ctl2, line width=0.8pt] table[x=t, y=x_pi] {images/p8_results/tikz/estable.csv}; \addlegendentry{PI}
  \addplot[tray, line width=0.8pt] table[x=t, y=x_pii] {images/p8_results/tikz/estable.csv}; \addlegendentry{PII (tesis)}
\end{axis}
\begin{axis}[pocho, at={(principal.south west)}, anchor=north west, yshift=-0.45cm, height=4.4cm,
    xmin=0, xmax=40, ymin=-0.2, ymax=3.8, xlabel={tiempo [\si{s}]}, ylabel={$u$ [\si{V}]}]
  \addplot[ctl2, line width=0.7pt] table[x=t, y=u_pi] {images/p8_results/tikz/estable.csv};
  \addplot[tray, line width=0.7pt] table[x=t, y=u_pii] {images/p8_results/tikz/estable.csv};
  \addplot[red!70!black, densely dotted, line width=0.8pt] coordinates {(0,3.5) (40,3.5)};
\end{axis}
\end{tikzpicture}

    \caption[Configuración estable]{Regulación en la configuración estable (repulsiva) con el PI y con el PII de la tesis, escalón de 1 a 2~cm. Arriba, posición; abajo, señal de control.}
    \label{fig:p8_estable}
\end{figure}

\FloatBarrier

Este capítulo deja un resultado: el PII estabiliza el imán en todas las pruebas, dentro de los límites del actuador, pero ningún escenario con atascamiento-deslizamiento termina en reposo sobre la referencia, porque el lazo oscila en un ciclo límite de unas centésimas de centímetro que el modelo lineal con zona muerta y fricción reproduce. Queda por explicar su origen y por saber si la segunda integración lo reduce, que es el tema del capítulo siguiente.
